\documentclass{article}
\usepackage{amsmath,amssymb,amsthm}
\usepackage{algorithm}
\usepackage{algorithmic}
\usepackage{xcolor}
\newcommand{\R}{\mathbb{R}}
\newcommand{\E}{\mathbb{E}}
\newcommand{\norm}[1]{\left\lVert #1\right\rVert}
\newcommand{\inner}[1]{\left\langle #1\right\rangle}
\newtheorem{theorem}{Theorem}
\newtheorem{lemma}{Lemma}
\newtheorem{assumption}{Assumption}
\newtheorem{remark}{Remark}

\begin{document}

\section*{Differentially Private Stochastic Gradient Descent (DP‑SGD)}

\section*{Mathematical Formulation}
Let $f:\R^{d}\to\R$ be the empirical risk
\[
f(w)=\frac{1}{n}\sum_{i=1}^{n}\ell(w;x_i),\qquad 
\ell(\cdot;\cdot)\;\text{convex and $L$‑smooth}.
\]
At iteration $t$ we draw a minibatch $B_t\subset\{1,\dots,n\}$ of size $b$ uniformly at random.
For each $i\in B_t$ compute the per‑example gradient
\[
g_i^{(t)}=\nabla \ell(w_t;x_i).
\]
Clip it to a fixed norm $C>0$:
\[
\tilde g_i^{(t)}=g_i^{(t)}\;\min\!\Big\{1,\frac{C}{\norm{g_i^{(t)}}}\Big\}.
\]
Average the clipped gradients and add isotropic Gaussian noise:
\[
\hat g_t=\frac{1}{b}\sum_{i\in B_t}\tilde g_i^{(t)}+ \xi_t,\qquad 
\xi_t\sim\mathcal N\!\big(0,\;\sigma^{2}C^{2}\,I_d\big).
\]
The DP‑SGD update is
\[
\boxed{\,w_{t+1}=w_t-\eta_t\,\hat g_t\,},
\]
where $\eta_t>0$ is the learning‑rate schedule.  
All random variables are defined on a common probability space; the noise $\xi_t$ is independent of the data and of previous iterates.

\section*{Pseudocode}
\begin{algorithm}[H]
\caption{DP‑SGD}
\begin{algorithmic}[1]
\STATE \textbf{Input:} data $\{x_i\}_{i=1}^n$, initial point $w_0$, learning rates $\{\eta_t\}_{t=0}^{T-1}$,
clipping norm $C$, noise scale $\sigma$, batch size $b$, total iterations $T$.
\FOR{$t=0$ \TO $T-1$}
    \STATE Sample minibatch $B_t\subset\{1,\dots,n\}$, $|B_t|=b$.
    \FOR{$i\in B_t$}
        \STATE $g_i^{(t)}\gets \nabla\ell(w_t;x_i)$
        \STATE $\tilde g_i^{(t)}\gets g_i^{(t)}\;\min\!\Big\{1,\frac{C}{\norm{g_i^{(t)}}}\Big\}$ \COMMENT{clipping}
    \ENDFOR
    \STATE $\bar g_t\gets \frac{1}{b}\sum_{i\in B_t}\tilde g_i^{(t)}$
    \STATE Sample $\xi_t\sim\mathcal N(0,\sigma^{2}C^{2}I_d)$
    \STATE $\hat g_t\gets \bar g_t+\xi_t$
    \STATE $w_{t+1}\gets w_t-\eta_t\hat g_t$
\ENDFOR
\STATE \textbf{Output:} $w_T$ (or the average $\bar w_T=\frac1T\sum_{t=0}^{T-1}w_t$)
\end{algorithmic}
\end{algorithm}

\section*{Dry Run Example}
Consider the scalar function $f(w)=(w-3)^2$, so $\nabla f(w)=2(w-3)$.  
Take $w_0=0$, learning rate $\eta=0.1$, clipping norm $C=10$ (no clipping occurs), noise variance $\sigma^2=0$ (for illustration), batch size $b=1$, and run $3$ iterations.

\begin{center}
\begin{tabular}{c|c|c|c}
$t$ & $g_t=\nabla f(w_t)$ & $w_{t+1}=w_t-\eta g_t$ & $f(w_{t+1})$\\ \hline
0 & $2(0-3)=-6$ & $0-0.1(-6)=0.6$ & $(0.6-3)^2=5.76$\\
1 & $2(0.6-3)=-4.8$ & $0.6-0.1(-4.8)=1.08$ & $(1.08-3)^2=3.6864$\\
2 & $2(1.08-3)=-3.84$ & $1.08-0.1(-3.84)=1.464$ & $(1.464-3)^2=2.3665$\\
\end{tabular}
\end{center}
The objective value decreases at each step, illustrating the descent property of the algorithm (in the noiseless case).

\newpage

\section*{Assumptions}
\begin{assumption}[Convexity and Smoothness]\label{ass:convex}
The loss $\ell(\cdot;x)$ is convex for every data point $x$ and $L$‑smooth, i.e.
\[
\ell(y;x)\le \ell(z;x)+\inner{\nabla\ell(z;x)}{y-z}+\frac{L}{2}\norm{y-z}^{2},
\qquad\forall y,z\in\R^{d}.
\]
Consequently $f$ is convex and $L$‑smooth.
\end{assumption}
\begin{assumption}[Bounded Clipped Gradients]\label{ass:clip}
After clipping, all gradients satisfy $\norm{\tilde g_i^{(t)}}\le C$ for all $i,t$.
\end{assumption}
\begin{assumption}[Unbiased Noisy Gradient]\label{ass:unbiased}
The noise $\xi_t$ is zero‑mean and independent of the data, thus
\[
\E[\hat g_t\mid w_t]=\E\!\Big[\frac{1}{b}\sum_{i\in B_t}\tilde g_i^{(t)}\;\Big|\;w_t\Big]=\nabla f(w_t).
\]
\end{assumption}
\begin{assumption}[Finite Variance]\label{ass:variance}
Conditioned on $w_t$, the variance of $\hat g_t$ is bounded:
\[
\E\!\big[\norm{\hat g_t-\nabla f(w_t)}^{2}\mid w_t\big]\le
\frac{C^{2}}{b}+ \sigma^{2}C^{2}d
=: \tau^{2}.
\]
(The first term comes from sampling without replacement, the second from Gaussian noise.)
\end{assumption}

\section*{Main Convergence Theorem}
\begin{theorem}\label{thm:conv}
Assume \ref{ass:convex}–\ref{ass:variance} and a constant stepsize $\eta_t\equiv\eta\le\frac{1}{L}$.
Let $w^{*}\in\arg\min_{w}f(w)$. Then the iterates of DP‑SGD satisfy
\[
\E\!\big[f(\bar w_T)-f(w^{*})\big]\;\le\;
\frac{\norm{w_0-w^{*}}^{2}}{2\eta T}
\;+\;\frac{\eta L}{2}\,\tau^{2},
\qquad
\bar w_T:=\frac{1}{T}\sum_{t=0}^{T-1}w_t .
\]
In particular, choosing $\eta=\Theta\!\big(\frac{1}{\sqrt{T}}\big)$ yields
\[
\E\!\big[f(\bar w_T)-f(w^{*})\big]=O\!\Big(\frac{1}{\sqrt{T}}\Big).
\]
\end{theorem}

\section*{Theoretical Properties}
\begin{itemize}
\item \textbf{Stability.} The clipping operator is non‑expansive, guaranteeing that the update never explodes regardless of the raw gradient magnitude.
\item \textbf{Hyperparameter Sensitivity.}
    \begin{itemize}
        \item Larger clipping norm $C$ reduces bias from clipping but increases noise variance $\sigma^{2}C^{2}d$.
        \item The stepsize $\eta$ trades off the optimization error $\frac{\norm{w_0-w^{*}}^{2}}{2\eta T}$ against the variance term $\frac{\eta L}{2}\tau^{2}$.
    \end{itemize}
\item \textbf{Comparison to Non‑Private SGD.}
    Setting $C\to\infty$ and $\sigma=0$ recovers the standard SGD bound
    $\E[f(\bar w_T)-f(w^{*})]\le\frac{\norm{w_0-w^{*}}^{2}}{2\eta T}+\frac{\eta L G^{2}}{2}$
    with $G$ a uniform gradient bound.
\item \textbf{Privacy.} The Gaussian mechanism with noise scale $\sigma C$ guarantees $(\varepsilon,\delta)$‑DP after $T$ compositions (via moments accountant or R\'enyi DP). The convergence theorem is independent of the privacy accounting; it only uses the statistical properties of the added noise.
\end{itemize}

\section*{Discussion}
The bound of Theorem~\ref{thm:conv} mirrors the classical SGD rate $O(1/\sqrt{T})$ while explicitly exposing the penalty incurred by differential privacy through the term $\sigma^{2}C^{2}d$. Choosing $C$ as small as permissible (subject to acceptable clipping bias) and calibrating $\sigma$ to the desired privacy budget directly controls the optimization error.

\newpage

\section*{Preliminaries (Definitions and Lemmas)}
\begin{lemma}[Smoothness Inequality]\label{lem:smooth}
If $f$ is $L$‑smooth then for any $w,u\in\R^{d}$
\[
f(u)\le f(w)+\inner{\nabla f(w)}{u-w}+\frac{L}{2}\norm{u-w}^{2}.
\]
\end{lemma}
\begin{proof}
Standard from the definition of $L$‑smoothness. \qedhere
\end{proof}

\begin{lemma}[Non‑expansiveness of Clipping]\label{lem:clip}
For any $a,b\in\R^{d}$ and clipping norm $C>0$,
\[
\norm{\operatorname{clip}_{C}(a)-\operatorname{clip}_{C}(b)}\le \norm{a-b},
\]
where $\operatorname{clip}_{C}(v)=v\min\{1,\frac{C}{\norm{v}}\}$.
\end{lemma}
\begin{proof}
Follows from the fact that the clipping operator is a projection onto the Euclidean ball of radius $C$. \qedhere
\end{proof}

\section*{Main Proof (Step‑by‑Step Derivation)}
We prove Theorem~\ref{thm:conv}. Throughout the proof, expectations are taken over all randomness up to iteration $t$ unless otherwise noted.

\bigskip
\noindent\textbf{Step 1. Distance Recursion.}  
From the update $w_{t+1}=w_t-\eta\hat g_t$,
\begin{align}
\norm{w_{t+1}-w^{*}}^{2}
&=\norm{w_t-w^{*}}^{2}
-2\eta\inner{\hat g_t}{w_t-w^{*}}
+\eta^{2}\norm{\hat g_t}^{2}. \tag{1}\label{eq:dist}
\end{align}
[Step 1] expands the squared norm without omission.

\bigskip
\noindent\textbf{Step 2. Take Conditional Expectation.}  
Condition on $w_t$ and use \ref{ass:unbiased} ($\E[\hat g_t\mid w_t]=\nabla f(w_t)$):
\begin{align}
\E\!\big[\norm{w_{t+1}-w^{*}}^{2}\mid w_t\big]
&=\norm{w_t-w^{*}}^{2}
-2\eta\inner{\nabla f(w_t)}{w_t-w^{*}}
+\eta^{2}\E\!\big[\norm{\hat g_t}^{2}\mid w_t\big]. \tag{2}\label{eq:cond}
\end{align}
[Step 2] moves the expectation inside and replaces $\hat g_t$ by its conditional mean in the inner product term.

\bigskip
\noindent\textbf{Step 3. Decompose the Second‑moment Term.}
\begin{align}
\E\!\big[\norm{\hat g_t}^{2}\mid w_t\big]
&=\norm{\nabla f(w_t)}^{2}
+\E\!\big[\norm{\hat g_t-\nabla f(w_t)}^{2}\mid w_t\big] \nonumber\\
&\le \norm{\nabla f(w_t)}^{2}+\tau^{2}, \tag{3}\label{eq:second}
\end{align}
where the inequality follows from Assumption~\ref{ass:variance}.  
[Step 3] uses the identity $\|a+b\|^{2}=\|a\|^{2}+2\langle a,b\rangle+\|b\|^{2}$ and the zero‑mean property of $\hat g_t-\nabla f(w_t)$.

\bigskip
\noindent\textbf{Step 4. Plug \eqref{eq:second} into \eqref{eq:cond}.}
\begin{align}
\E\!\big[\norm{w_{t+1}-w^{*}}^{2}\mid w_t\big]
&\le \norm{w_t-w^{*}}^{2}
-2\eta\inner{\nabla f(w_t)}{w_t-w^{*}}
+\eta^{2}\norm{\nabla f(w_t)}^{2}
+\eta^{2}\tau^{2}. \tag{4}\label{eq:plug}
\end{align}
[Step 4] is a direct substitution.

\bigskip
\noindent\textbf{Step 5. Relate the Inner Product to Function Values.}  
By convexity of $f$,
\[
f(w_t)-f(w^{*})\le \inner{\nabla f(w_t)}{w_t-w^{*}}. \tag{5}\label{eq:conv}
\]
[Step 5] is the standard first‑order condition for convex functions.

\bigskip
\noindent\textbf{Step 6. Apply Smoothness.}  
Lemma~\ref{lem:smooth} yields
\[
f(w_{t+1})\le f(w_t)+\inner{\nabla f(w_t)}{w_{t+1}-w_t}
+\frac{L}{2}\norm{w_{t+1}-w_t}^{2}. \tag{6}\label{eq:smooth}
\]
Since $w_{t+1}-w_t=-\eta\hat g_t$, we obtain
\[
f(w_{t+1})\le f(w_t)-\eta\inner{\nabla f(w_t)}{\hat g_t}
+\frac{L\eta^{2}}{2}\norm{\hat g_t}^{2}. \tag{7}\label{eq:smooth2}
\]

\bigskip
\noindent\textbf{Step 7. Take Conditional Expectation of \eqref{eq:smooth2}.}
\begin{align}
\E\!\big[f(w_{t+1})\mid w_t\big]
&\le f(w_t)-\eta\inner{\nabla f(w_t)}{\nabla f(w_t)}
+\frac{L\eta^{2}}{2}\E\!\big[\norm{\hat g_t}^{2}\mid w_t\big] \nonumber\\
&= f(w_t)-\eta\norm{\nabla f(w_t)}^{2}
+\frac{L\eta^{2}}{2}\big(\norm{\nabla f(w_t)}^{2}+\tau^{2}\big) \tag{8}\label{eq:exp_smooth}
\end{align}
where we used \ref{ass:unbiased} and \eqref{eq:second}.  
[Step 7] rearranges terms and substitutes the bound on the second moment.

\bigskip
\noindent\textbf{Step 8. Combine \eqref{eq:plug} and \eqref{eq:exp_smooth}.}  
Subtract $2\eta\big(f(w_t)-f(w^{*})\big)$ from both sides of \eqref{eq:plug} (using \eqref{eq:conv}) and then add \eqref{eq:exp_smooth}:
\begin{align}
\E\!\big[\norm{w_{t+1}-w^{*}}^{2}\mid w_t\big]
&\le \norm{w_t-w^{*}}^{2}
-2\eta\big(f(w_t)-f(w^{*})\big)
+\eta^{2}\tau^{2}
\tag{9}\label{eq:rec1}\\
\E\!\big[f(w_{t+1})\mid w_t\big]
&\le f(w_t)-\eta\big(1-\tfrac{L\eta}{2}\big)\norm{\nabla f(w_t)}^{2}
+\frac{L\eta^{2}}{2}\tau^{2}. \tag{10}\label{eq:rec2}
\end{align}
[Step 8] isolates the optimality gap in \eqref{eq:rec1}.

\bigskip
\noindent\textbf{Step 9. Choose $\eta\le\frac{1}{L}$}.  
Then $1-\frac{L\eta}{2}\ge\frac{1}{2}$, and the non‑negative term $-\eta\big(1-\frac{L\eta}{2}\big)\norm{\nabla f(w_t)}^{2}$ can be dropped to obtain a simpler recursion:
\[
\E\!\big[f(w_{t+1})\mid w_t\big]\le f(w_t)+\frac{L\eta^{2}}{2}\tau^{2}. \tag{11}\label{eq:simple}
\]

\bigskip
\noindent\textbf{Step 10. Unroll the Distance Recursion.}  
Taking total expectation in \eqref{eq:rec1} and summing for $t=0,\dots,T-1$ gives
\begin{align}
\E\!\big[\norm{w_T-w^{*}}^{2}\big]
&\le \norm{w_0-w^{*}}^{2}
-2\eta\sum_{t=0}^{T-1}\E\!\big[f(w_t)-f(w^{*})\big]
+T\eta^{2}\tau^{2}. \tag{12}\label{eq:sum}
\end{align}
Rearranging,
\[
\frac{1}{T}\sum_{t=0}^{T-1}\E\!\big[f(w_t)-f(w^{*})\big]
\le \frac{\norm{w_0-w^{*}}^{2}}{2\eta T}
+\frac{\eta\tau^{2}}{2}. \tag{13}\label{eq:avg_gap}
\]

\bigskip
\noindent\textbf{Step 11. From Iterate Average to Final Output.}  
By convexity,
\[
f(\bar w_T)-f(w^{*})\le\frac{1}{T}\sum_{t=0}^{T-1}\big(f(w_t)-f(w^{*})\big).
\]
Taking expectation and using \eqref{eq:avg_gap} yields
\[
\E\!\big[f(\bar w_T)-f(w^{*})\big]
\le\frac{\norm{w_0-w^{*}}^{2}}{2\eta T}
+\frac{\eta\tau^{2}}{2}. \tag{14}\label{eq:final}
\]
Since $\tau^{2}=C^{2}\big(\frac{1}{b}+\sigma^{2}d\big)$, the theorem follows. \qed

\section*{Rate Derivation (With Constants)}
From \eqref{eq:final} set $\eta=\frac{1}{L}\sqrt{\frac{\norm{w_0-w^{*}}^{2}}{T\tau^{2}}}$ (which respects $\eta\le 1/L$). Then
\[
\E\!\big[f(\bar w_T)-f(w^{*})\big]
\le \frac{L}{\sqrt{T}}\norm{w_0-w^{*}}\tau
= O\!\Big(\frac{1}{\sqrt{T}}\Big).
\]
Explicitly,
\[
\E\!\big[f(\bar w_T)-f(w^{*})\big]
\le
\frac{L\,\norm{w_0-w^{*}}\,C}{\sqrt{T}}\sqrt{\frac{1}{b}+\sigma^{2}d}.
\]

\section*{Special Cases}
\begin{itemize}
\item \textbf{No privacy noise} ($\sigma=0$). The bound reduces to the classic SGD rate with variance term $C^{2}/b$ stemming only from minibatch sampling.
\item \textbf{Full‑batch ($b=n$) and no clipping} ($C\ge\max_i\norm{\nabla\ell(w;x_i)}$). Then $\tau^{2}= \sigma^{2}C^{2}d$ and the error is solely due to the added Gaussian noise.
\item \textbf{Large $T$ limit}. The dominant term is $O(1/\sqrt{T})$; the privacy‑induced variance appears only as a constant factor inside the square root.
\end{itemize}

\end{document}